This draft supports the recoverable-resilience framing. The empirical data reveal rainfall-aligned mobility disruptions, nonuniform speed-deficit burden, OD functional dependence, and heterogeneous natural recovery. The optimization model then estimates the counterfactual portion of this loss that is recoverable under a standardized management regime.

The learning analysis extracts two linked laws from the optimizer-derived recovery-value field. First, local marginal recovery value is activated when persistent future deficit overlaps with OD exposure and feasible intervention efficiency. Second, finite-budget allocation must be residual: after each deployment, the remaining value depends on the current unrecovered loss, remaining budget, remaining time, and diminishing returns. At the event level, decision-criticality is governed by the top tail of marginal recovery values rather than disruption magnitude alone.

The current results should be read as regime-conditional evidence, not as a final universal law. The next strongest extensions are wider scenario-specific LP closure, intervention-parameter sensitivity, and a graph-based surrogate for larger-scale law extraction.
